% Source-derived chapter 11: Local intersection number and trace formula
% Original source: higher-siegel-weil-unitary-nonsingular-terms
\chapter{Local intersection number and trace formula}
\label{higher-siegel-weil-unitary-nonsingular-terms-chapter-11}
\source{higher-siegel-weil-unitary-nonsingular-terms}

\begin{remark}[Source section]
\label{higher-siegel-weil-unitary-nonsingular-terms-chapter-11-source}
\source{higher-siegel-weil-unitary-nonsingular-terms}
\source{higher-siegel-weil-unitary-nonsingular-terms}
This chapter reproduces the corresponding section of the retrieved
LaTeX source, including its definitions, statements, examples, and
proofs. It has not yet been translated into Lean.
\notready
\end{remark}

% The source section title was: Local intersection number and trace formula
\label{s:int trace}
\source{higher-siegel-weil-unitary-nonsingular-terms}

\subsection{Local nature of the intersection problem}Recall from Proposition \ref{prop: KR cycle proper} that $\cZ^{r}_{\cE}(a)$ only depends on the Hermitian torsion sheaf $\cQ=\coker(a)$. In this subsection we show that the $0$-cycle class $[\cZ^{r}_{\cE}(a)]$ also only depends on $\cQ$. 

Recall the stacks $\Herm_{2d}=\Herm_{2d}(X'/X)$ and $\Lagr_{2d}$ from \S \ref{sec: Herm}. We have a self-correspondence $\Hk^{r}_{\Lagr_{2d}}$ of $\Lagr_{2d}$ over $\Herm_{2d}$: it classifies 
$(\cQ,h,\{\cL_{i}\}_{0\le i\le r})$ where $(\cQ,h)\in \Herm_{2d}$, $\cL_{i}\subset \cQ$ are Lagrangian subsheaves such that $\cL_{i}/(\cL_{i}\cap \cL_{i-1})$ has length one for $1\le i\le r$. 
%\begin{equation}\label{corr Herm}
%\xymatrix{& \Hk^{r}_{\Lagr_{2d}} \ar[dl]\ar[dr] \\
%\Lagr_{2d}\ar[dr]_{\ph} && \Lagr_{2d}\ar[dl]^{\ph}\\
%& \Herm_{2d}}
%\end{equation}
Define the local version $\Sht^{r}_{\Lagr_{2d}}$ of $\Sht^{r}_{\cM_{d}}$ by the Cartesian diagram
\begin{equation}\label{Sht Lag}
\source{higher-siegel-weil-unitary-nonsingular-terms}
\xymatrix{\Sht^{r}_{\Lagr_{2d}}\ar[d] \ar[r] & \Hk^{r}_{\Lagr_{2d}}\ar[d]^{(\pr_{0},\pr_{r})}\\
\Lagr_{2d}\ar[r]^-{(\Id,\Frob)} & \Lagr_{2d}\times\Lagr_{2d}}
\end{equation}
We have a decomposition into open-closed substacks
\begin{equation}
\Sht^{r}_{\Lagr_{2d}}=\coprod_{(\cQ,h)\in\Herm_{2d}(k)} \cZ^{r}_{\cQ}.
\end{equation}

\begin{lemma}
\source{higher-siegel-weil-unitary-nonsingular-terms}
\notready
The stack $\Hk^{1}_{\Lagr_{2d}}$ is smooth of dimension zero.
\end{lemma}
\begin{proof}
We may identify $\Hk^{1}_{\Lagr_{2d}}$ with the moduli stack of $(0\subset \cL'\subset \cL\subset \cQ, h)$ where $(\cL\subset\cQ,h)\in\Lagr_{2d}$ and $\cL/\cL'$ has length one. Under the local chart for $\Herm_{2d}$ described in Lemma \ref{l:Herm local}, $\Hk^{1}_{\Lagr_{2d}}$ becomes $[\frp/P]$, where $P\subset \Og_{2d}=\Og(V)$ is the parabolic subalgebra stabilizing a pair of subspaces  $L'\subset L$ with $L$ Lagrangian and $\dim L'=d-1$. This local description implies that $\Hk^{1}_{\Lagr_{2d}}$ is smooth of dimension zero.
\end{proof}

Rewriting $\Sht^{r}_{\Lagr_{2d}}$ as the fibered product (cf. Definition \ref{def:Phi})
\begin{equation}\label{Sht Lagr}
\source{higher-siegel-weil-unitary-nonsingular-terms}
\xymatrix{\Sht^{r}_{\Lagr_{2d}}\ar[r]\ar[d] & (\Hk^{1}_{\Lagr_{2d}})^{r} \times \Lagr_{2d} \ar[d]^{(\pr_{0},\pr_{1})^{r} \times \Delta}\\
(\Lagr_{2d})^{r+1}\ar[r]^{\Phi_{\Lagr_{2d}}} & (\Lagr_{2d})^{2r+2}}
\end{equation}
we define a $0$-cycle class
\begin{equation}
[\Sht^{r}_{\Lagr_{2d}}]:=\Phi_{\Lagr_{2d}}^{!}[(\Hk^{1}_{\Lagr_{2d}})^{r} \times \Lagr_{2d}]\in \Ch_{0}(\Sht^{r}_{\Lagr_{2d}}).
\end{equation}
Restricting to $\cZ^{r}_{\cQ}$ we get
\begin{equation}
[\cZ^{r}_{\cQ}]:=[\Sht^{r}_{\Lagr_{2d}}]|_{\cZ^{r}_{\cQ}}\in \Ch_{0}(\cZ^{r}_{\cQ}).
\end{equation}

Recall the maps $g_{\cM}: \cM_{d}\to \Lagr_{2d}$ and $g:\cA_{d}\to \Herm_{2d}$ defined in \S\ref{ss:vb to tor}. We also have a map  $g_{\Hk}: \Hk^{r}_{\cM_{d}}\to \Hk^{r}_{\Lagr_{2d}}$ sending $(\{x'_{i}\} , \{\cE\xr{t_{i}}\cF_{i}\})$ to $(\cQ=\coker(a), h_{\cQ}, \{\coker(t_{i})\})$ ($a$ is the induced Hermitian map on $\cE$).   The maps $g_{\Hk}, g_{\cM}$ and $g$ exhibit the diagram \eqref{def ShtMd} as the pullback of the diagram \eqref{Sht Lag} via the base change $g: \cA_{d}\to\Herm_{2d}$.  In particular we have a natural map
\begin{equation}
g_{\Sht}: \Sht^{r}_{\cM_{d}}\to \Sht^{r}_{\Lagr_{2d}}.
\end{equation}
For fixed $(\cE,a)\in \cA_{d}(k)$ with image $\cQ=\coker(a)\in \Herm_{2d}(k)$, $g_{\Sht}$ restricts to an isomorphism to the open-closed subschemes
\begin{equation}\label{gSht}
\source{higher-siegel-weil-unitary-nonsingular-terms}
g_{\Sht}|_{\cZ^{r}_{\cE}(a)}: \cZ^{r}_{\cE}(a)\isom \cZ^{r}_{\cQ}.
\end{equation}


% Since $g\co\cA_{d}\to\Herm_{2d}$ is smooth, Proposition \ref{p:base change} implies:
\begin{prop}
\source{higher-siegel-weil-unitary-nonsingular-terms}
\notready We have an equality
\begin{equation}\label{ShtM ShtLag}
\source{higher-siegel-weil-unitary-nonsingular-terms}
[\Sht^{r}_{\cM_{d}}]=g_{\Sht}^{*}[\Sht^{r}_{\Lagr_{2d}}]\in \Ch_{0}(\Sht^{r}_{\cM_{d}}).
\end{equation}
\end{prop} 
\begin{proof}
Apply Proposition \ref{p:base change} to the diagram \eqref{Sht Lagr}, the fundamental class $\z=[(\Hk^{1}_{\Lagr_{2d}})^{r} \times \Lagr_{2d}]$ and the base change map $u=g: \cA_{d}\to \Herm_{2d}$. By Proposition \ref{prop:Ad Herm smooth}, $g$ is smooth. We then have
\begin{equation}
\Phi_{\cM_{d}}^{!}g_{\Hk}^{*}[(\Hk^{1}_{\Lagr_{2d}})^{r} \times \Lagr_{2d}]=g^{*}_{\Sht}\Phi_{\Lagr_{2d}}^{!}[(\Hk^{1}_{\Lagr_{2d}})^{r} \times \Lagr_{2d}]\in \Ch_{0}(\Sht^{r}_{\cM_{d}}).
\end{equation}
Since $g_{\Hk}$ is smooth,  $g_{\Hk}^{*}[(\Hk^{1}_{\Lagr_{2d}})^{r} \times \Lagr_{2d}]=[(\Hk^{1}_{\cM_{d}})^{r} \times \Hk^{1}_{\cM_{d}}]$. The above equality then becomes \eqref{ShtM ShtLag}.
\end{proof}


Combined with Theorem \ref{th:compare KR}, we get a local description of the cycle class $[\cZ^{r}_{\cE}(a)]$:
\begin{cor}
\source{higher-siegel-weil-unitary-nonsingular-terms}
\notready\label{c:Zloc}For any $(\cE,a)\in \cA_{d}(k)$ with image $\cQ=\coker(a)\in \Herm_{2d}(k)$, $[\cZ^{r}_{\cE}(a)]$ is the same as $[\cZ^{r}_{\cQ}]$ under the isomorphism \eqref{gSht}. In particular,
\source{higher-siegel-weil-unitary-nonsingular-terms}
\begin{equation}
\deg[\cZ^{r}_{\cE}(a)]=\deg[\cZ^{r}_{\cQ}].
\end{equation}
\end{cor}



\subsection{Sheaves on $\Herm_{2d}$}
To describe the direct image complex $Rf_{*}\Qlbar$ on $\cA_{d}$, by the Cartesian diagram \eqref{eq: hitchin cartesian}, we first need to understand $R(\upsilon_{2d})_*\Qlbar$ on $\Herm_{2d}$.

\begin{lemma}
\source{higher-siegel-weil-unitary-nonsingular-terms}
\notready\label{l:Lag push}
\source{higher-siegel-weil-unitary-nonsingular-terms}
The perverse sheaf $R(\upsilon_{2d})_{*}\Qlbar$ on $\Herm_{2d}$ is canonically isomorphic to $(\Spr^{\Herm}_{2d})^{S_{d}}$ (see Proposition \ref{p:SH Spr}(2)). Here the $S_{d}$-action on $\Spr^{\Herm}_{2d}$ is the restriction of the Springer $W_{d}$-action.
\end{lemma}
\begin{proof} We have a Cartesian diagram
\begin{equation}
\xymatrix{\wt\Herm_{2d}\ar[d]^{\l_{2d}} \ar[r]^{\wt\e_{d}} & \wt\Coh_{d}(X') \ar[d]^{\pi^{\Coh}_{X',d}} \\
\Lagr_{2d} \ar[r]^{\e'_{d}} & \Coh_{d}(X')
}
\end{equation}
where $\e'_{d}$ sends $(\cQ,h_{\cQ},\cL)$ to $\cL$ and $\wt\e_{d}$ sends $(\cQ_{1}\subset \cdots\subset\cQ_{d}\subset \cdots\subset \cQ,h)$ to $(\cQ_{1}\subset \cdots\subset\cQ_{d})$. By proper base change
\begin{equation}
R\l_{2d*}\Qlbar\cong \e'^{*}_{d}\Spr_{d, X'}.
\end{equation}
In particular, $R\l_{2d*}\Qlbar$ carries an action of $S_{d}$. Moreover, the induced $S_{d}$-action  on  $\Spr^{\Herm}_{2d}\cong R(\upsilon_{2d})_{*}R\l_{2d*}\Qlbar$ is  the restriction  of the Springer $W_{d}$-action to $S_{d}$:  this can be easily checked over $\Herm^{\c}_{2d}$, and then the statement holds over $\Herm_{2d}$ since $\Spr^{\Herm}_{2d}$ is the middle extension from its restriction to $\Herm^{\c}_{2d}$ by Proposition \ref{p:SH Spr}(2). Since  $(R\l_{2d*}\Qlbar)^{S_{d}}\cong \e'^{*}_{d}(\Spr_{d,X'})^{S_{d}}\cong \Qlbar$, we conclude that $(\Spr^{\Herm}_{2d})^{S_{d}}\cong R(\upsilon_{2d})_{*}(R\l_{2d*}\Qlbar)^{S_{d}}\cong R(\upsilon_{2d})_{*}\Qlbar$, as desired.
\end{proof}

It is an elementary exercise to see that $\Ind_{S_{d}}^{W_{d}}\one$ decomposes into irreducible representations
\begin{equation}\label{rho i}
\source{higher-siegel-weil-unitary-nonsingular-terms}
\Ind_{S_{d}}^{W_{d}}\one=\bigoplus_{i=0}^{d}\r_{i}
\end{equation}
where $\r_{i}=\Ind^{W_{d}}_{(\Z/2\Z)^{d}\rtimes (S_{i}\times S_{d-i})}(\chi_{i}\bt \one)$, and $\chi_{i}: (\Z/2\Z)^{d}\rtimes (S_{i}\times S_{d-i})\to \{\pm1\}$ is the character that is nontrivial on the first $i$ factors of $(\Z/2\Z)^{i}$, trivial on the rest and trivial on $S_{i}\times S_{d-i}$.  The decomposition also shows up in \cite[\S8.1.1]{YZ}.

Recall the notation $\Spr^{\Herm}_{2d}[\r]$ from Definition \ref{d:Spr isotypic}.
\begin{cor}
\source{higher-siegel-weil-unitary-nonsingular-terms}
\notready\label{c:Lag decomp} There is a canonical decomposition 
\source{higher-siegel-weil-unitary-nonsingular-terms}
\begin{equation}\label{decomp ph}
\source{higher-siegel-weil-unitary-nonsingular-terms}
R(\upsilon_{2d})_{*}\Qlbar\cong\bigoplus_{i=0}^{d}\Spr^{\Herm}_{2d}[\r_{i}].
\end{equation}
\end{cor}
\begin{proof}
By Lemma \ref{l:Lag push} and Frobenius reciprocity,  we have
\begin{equation}
R(\upsilon_{2d})_{*}\Qlbar\cong \Hom_{W_{d}}(\Ind_{S_{d}}^{W_{d}}\one, \Spr^{\Herm}_{2d})=\Spr^{\Herm}_{2d}[\Ind_{S_{d}}^{W_{d}}\one].
\end{equation}
The desired decomposition then follows from \eqref{rho i}.
\end{proof}


\begin{defn}
\source{higher-siegel-weil-unitary-nonsingular-terms}
\notready\label{def:Kd}
\source{higher-siegel-weil-unitary-nonsingular-terms}
Define the graded perverse sheaf on $\Herm_{2d}(X'/X)$ 
\begin{equation}
\cK_{d}^{\mrm{Int}}(T) := \bigoplus_{i=0}^d \Spr^{\Herm}_{2d}[\r_{i}] T^{i}.
\end{equation}
\end{defn}

%Using the local chart for $\Herm_{2d}$ in Lemma \ref{l:Herm local}, the sequence \eqref{factor pi Herm} can be identified \'etale locally with
%\begin{equation}
%\pi_{2d}: [\wt \so_{2d}/\Og_{2d}] \xr{\l_{2d}} [\wt\so^{(d)}_{2d}/\Og_{2d}]\xr{\upsilon_{2d}} [\so_{2d}/\Og_{2d}]
%\end{equation}
%Here $\wt\so^{(d)}_{2d}=\{(A,L)|A\in \so_{2d}, L\subset V \mbox{ is a Lagrangian},  A(L)\subset L\}$. By The complex $R\l_{2d*}\Qlbar$ carries the action of $S_{d}\subset W_{d}$, such that the induced $S_{d}$-action on $R\pi_{2d*}\Qlbar$ is the restriction  of the Springer $W_{d}$-action to $S_{d}$. 


%To $(Q, h_Q) \in \Herm_d$, we can associate $\supp(Q)$ \tony{this isn't great terminology -- it's not actually the support in the literal sense} which descends to $X$. This defines a map $\Herm_d \rightarrow X_d$. Similarly, to $(Q, h_Q, L) \in \Lagr_d$ we can associate $\supp(L) \in X_d'$. 

%We have a map $s^{\Lagr}_{2d} \co \Lagr_{2d}\to X'_{d}$ sending $(\cQ,h_{\cQ},\cL)$ to the divisor of $\cL$. Recall also the map $s^{\Herm}_{2d}: \Herm_{2d}\to X_{d}$ sending $(\cQ,h_{\cQ})$ to the descent of the divisor of $\cQ$ to $X$.
%
%\begin{lemma} The diagram 
%\begin{equation}
%\xymatrix{\Lagr_{2d}\ar[d]^{\upsilon_{2d}} \ar[r]^{s^{\Lagr}_{2d}} & X_d' \ar[d]^{\nu_{d}} \\
%\Herm_{2d} \ar[r]^{s^{\Herm}_{2d}} & X_d}
%\end{equation}
%commutes and is Cartesian on $X_d^{\circ}$. 
%\end{lemma}
%\begin{proof}
%
%\end{proof}
%
%Note that $(X_d)^{\circ}$ has an \'{e}tale cover by $((X')^d)^{\circ}$, with Galois group $\Gamma_d := \{\pm 1 \}^d \rtimes S_d$. 
%
%We know from \cite[\S 8.1]{YZ} that 
%\[
%R\pi_{d!} \ol{\Q}_{\ell} \cong \bigoplus_{i=0}^d K_i, \quad \rank K_i = \binom{d}{i}
%\]
%where $K_i$ is the intermediate extension of the representation of $\Gamma_d$ obtained as follows: the character of $\{\pm 1 \}^d$ which is non-trivial on the first $i$ factors and trivial on the remaining $(d-i)$ factors extends to a character of $\{\pm 1 \}^d \rtimes S_i \times S_{d-i}$; we induce this to $\Gamma_d$. 

The fundamental class of the self-correspondence $\Hk^{1}_{\Lagr_{2d}}$ of $\Lagr_{2d}$ is viewed as a cohomological correspondence of the constant sheaf on $\Lagr_{2d}$ with itself. It induces an endomorphism (see notation from \cite[A.4.1]{YZ})
\begin{equation}
(\upsilon_{2d})_{!}[\Hk^{1}_{\Lagr_{2d}}]:R(\upsilon_{2d})_{*}\Qlbar\to R(\upsilon_{2d})_{*}\Qlbar.
\end{equation}


\begin{prop}
\source{higher-siegel-weil-unitary-nonsingular-terms}
\notready\label{p:action Hk1} The action of $(\upsilon_{2d})_{!}[\Hk^{1}_{\Lagr_{2d}}]$ on $R(\upsilon_{2d})_{*}\Qlbar$ preserves the decomposition \eqref{decomp ph},  and it acts on $\Spr^{\Herm}_{2d}[\r_{i}]$ by multiplication by $(d-2i)$.
\source{higher-siegel-weil-unitary-nonsingular-terms}
\end{prop}
\begin{proof}By Proposition \ref{p:SH Spr}(2),  $\Spr^{\Herm}_{2d}[\r_{i}]$ is the middle extension from its restriction to $\Herm^{\c}_{2d}$, it suffices to prove the same statement on $\Herm_{2d}^{\c}$. Let $I_d' \subset X_d' \times X'$ be the universal divisor $\{(D \in X_d', y \in X') \co y \in D\}$. The maps $\pr(D,y) = D$ and $q(D,y) = D-y+\sigma(y)$ define the incidence correspondence as in \cite[proof of Proposition 8.3]{YZ},
\begin{equation}\label{eq: incidence correspondence}
\source{higher-siegel-weil-unitary-nonsingular-terms}
\begin{tikzcd}
& I_d' \ar[dl, "\pr"'] \ar[dr, "q"] \\
X_d'  & & X_d' 
\end{tikzcd}
\end{equation}
Now over $\Herm_{2d}^{\c}$, the map $\Herm_{2d}^{\c}\to X^{\c}_{d}$ is smooth, and $\Hk^{1}_{\Lagr_{2d}}$ is the pullback of \eqref{eq: incidence correspondence}. We therefore reduce to checking the statement for the action of $[I'_{d}]$ on the direct image sheaf of $\nu_{d}: X'_{d}\to X_{d}$, which is done in the proof of \cite[Proposition 8.3]{YZ}.
\end{proof}

\subsection{Lefschetz trace formula}
We shall give a slight generalization of the Lefschetz trace formula \cite[Proposition A.12]{YZ} expressing the intersection number of a cycle with the graph of Frobenius as a trace. Instead of the graph of Frobenius, we need to intersect along $\Phi_{M}:M^{r+1}\to M^{2r+2}$. Consider the following situation:
\begin{itemize} 
\item Let $S$ be an algebraic stack locally of finite type over $k=\F_q$. Assume $S$ can be stratified by locally closed substacks that are global quotients.
\item Let $M$ be a smooth equidimensional stack over $k=\F_q$ of dimension $N$ with a proper representable map $f: M\to S$. 
\item For $1\le i\le r$,  let $(\pr^{i}_{0},\pr^{i}_{1}): C_{i}\to M\times_{S}M$ be a self-correspondence of $M$ over $S$. Assume $\pr^{i}_{0}$ is proper and representable.
\end{itemize}
Form the Cartesian diagram
\begin{equation}\label{ShtC}
\source{higher-siegel-weil-unitary-nonsingular-terms}
\xymatrix{\Sht_{C} \ar[r]\ar[d] & (\prod_{i=1}^{r}C_{i} ) \times M\ar[d]^{(\pr^{i}_{0}, \pr^{i}_{1})_{1\le i\le r}}\\
M^{r+1}\ar[r]^{\Phi_{M}} & M^{2r+2}\,.}
\end{equation}
Then $\Sht_{C}$ decomposes as
\begin{equation}
\Sht_{C}=\coprod_{s\in S(k)}\Sht_{C}(s)
\end{equation}
where $\Sht_C(s)$ is the fibered product of $\Sht_C \rightarrow S(k)$ against $s \in S(k)$. For $s \in S(k)$, we write $\Sht_{C}|_s$ for the fibered product 
\[
\begin{tikzcd}
\Sht_C|_s \ar[r] \ar[d] & \Sht_C \ar[d] \\
\Spec k \ar[r, "s"] & S(k)
\end{tikzcd}
\]
Then $\Sht_C|_s \rightarrow \Sht_C(s)$ is a torsor for $\Aut(s)$, and in particular a finite \'{e}tale cover. 

Suppose we are given cycle classes
\begin{equation}
\z_{i}\in \Ch_{N}(C_{i}), \quad 1\le i\le r.
\end{equation}
The cycle class  $\cla(\z_{i})\in \hBM{2N}{C_{i},\Qlbar(-N)}$ is viewed as a cohomological correspondence between the constant sheaf on $M$ and itself. Therefore it induces an endomorphism of $Rf_{!}\Qlbar$ which we denote by $f_{!}\cla(\z_{i})$.

\begin{prop}
\source{higher-siegel-weil-unitary-nonsingular-terms}
\notready\label{p:Lef} For each $s\in S(k)$ we have
\source{higher-siegel-weil-unitary-nonsingular-terms}
\begin{equation*}
\deg((\Phi_{M}^{!}(\z_{1}\times\cdots\times\z_{r} \times [M] ))|_{\Sht_{C}|_s})=\Tr(f_{!}\cla(\z_{1})\c\cdots\c f_{!}\cla(\z_{r})\c\Frob_{s}, (Rf_{!}\Qlbar)_{s}).
\end{equation*}
\end{prop}
\begin{proof} We first prove the formula when $S$ is a scheme of finite type.  In this case $M,C_{i}$ are also schemes of finite type over $k$. Let $C=C_{1}\times_{M}\times\cdots\times_{M}C_{r}$ be the composition correspondence, with maps $\pr_{i}: C\to M$ for $0\le i\le r$.  Consider the diagram where all squares are Cartesian
\begin{equation}
\xymatrix{\Sht_{C}\ar[d] \ar[r] & C \ar[r]\ar[d]^{(\pr_{0},\cdots, \pr_{r}, \pr_r)} & (\prod_{i=1}^{r}C_{i} ) \times M\ar[d]^{\prod (\pr^{i}_{0}, \pr^{i}_{1}) \times \Delta}\\
M^{r+1} \ar[d]_{\pr_{0}} \ar[r]^{\Phi_{1}} & M^{r+2}\ar[d]^{(\pr_{0}, \pr_{r+1})} \ar[r]^{\Phi_{2}} & M^{2r+2}\\
M\ar[r]_-{(\Id,\Frob_{M})} & M\times M}
\end{equation}
Here
\begin{align*}
\Phi_{1}(\xi_{0},\cdots,  \xi_{r-1}, \xi_r) & =(\xi_{0},\cdots, \xi_r, \Frob_{M}(\xi_{0})),\\
\Phi_{2}(\xi_{0},\cdots, \xi_r,\xi_{r+1}) & =(\xi_{0},\xi_{1},\xi_{1},\cdots,\xi_{r},\xi_{r}, \xi_{r+1}).
\end{align*}
We have $\Phi_{M}=\Phi_{2}\c\Phi_{1}$. Let $\z=\Phi_{2}^{!}(\z_{1}\times\z_{2}\times\cdots\times\z_{r} \times [M])\in \Ch_{N}(C)$. 

On the one hand, by the transitivity of the Gysin maps, 
\begin{equation}\label{Phizz}
\source{higher-siegel-weil-unitary-nonsingular-terms}
\Phi^{!}_{M}(\z_{1}\times\cdots\times\z_{r} \times [M] )=\Phi_{1}^{!}(\z)=(\Id,\Frob_{M})^{!}(\z).
\end{equation}
Applying the Lefschetz trace formula \cite[Proposition  A.12]{YZ}, we get
\begin{equation}\label{deg tr z}
\source{higher-siegel-weil-unitary-nonsingular-terms}
\deg((\Id,\Frob_{M})^{!}(\z)|_{\Sht_{C}|_s})=\Tr(f_{!}\cla(\z)\c \Frob, (Rf_{!}\Qlbar)_{s}).
\end{equation}
One the other hand, by a diagram chase, we see that $\cla(\z)$ is the composition of the cohomological correspondences $\cla(\z_{i})$ ($1\le i\le r$), hence $f_{!}\cla(\z)\in \End(Rf_{!}\Qlbar)$ is the composition of $f_{!}\cla(\z_{1})\c f_{!}\cla(\z_{2})\c\cdots\c f_{!}\cla(\z_{r})$. Combining this fact with \eqref{Phizz} and \eqref{deg tr z} we get the desired formula.

Now consider the general case where $S$ is a stack locally of finite type over $k$ and we aim to prove the formula for $s\in S(k)$. We claim that there exists a scheme $S'$ of finite type over $k$ and a smooth map $u: S'\to S$ such that $u(S'(k))$ contains $s$.  Indeed, pick any smooth map $u_{1}: S_{1}\to S$ with $S_{1}$ a scheme of finite type over $k$, such that  $s$ is contained in the image of $u$. Let $s_{1}\in S_{1}(\F_{q^{m}})$ be a point that maps to $s$. Let $(S_{1}/S)^{m}$ be the $m$-fold fibered product of $S_1$ over $S$, based changed to $\ol{k}$. We equip  $(S_{1}/S)^{m}$ with the $\Frob$-descent datum given by $(x_{1},\cdots, x_{m})\mapsto (\Frob(x_{m}), \Frob(x_{1}),\cdots, \Frob(x_{m-1}))$. This gives a descent of $(S_{1}/S)^{m}$ to a scheme $S'$ over $k$ equipped with a map $u:S'\to S$ which is still smooth since $u_{1}$ is. Now $s_{1}$ gives rise to a $k$-point $s'=(s_{1},\Frob(s_{1}),\cdots, \Frob^{m-1}(s_{1}))\in S'(k)$ such that $u(s')=s$. 

Let $M'=M\times_{S}S'$, $C'_{i}=C_{i}\times_{S}S'$ and let $u_{C_{i}}:C'_{i}\to C_{i}$ be the projection. Define $\Sht'_{C}$ using the analog of the diagram \eqref{ShtC} with $M$ and $C_{i}$ replaced by $M'$ and $C'_{i}$. Then $\Sht'_{C}\cong \Sht_{C}\times_{S(k)}S'(k)$. For $s'\in S'(k)$ such that $u(s')=s$, we get an isomorphism $\Sht'_{C}|_{s'}\isom \Sht_{C}|_{s}$. Let $\z'_{i}=u_{C_{i}}^{*}\z_{i}$. Now we apply Proposition \ref{p:base change} to the diagram \eqref{ShtC} along the base change map $u:S'\to S$ to get
\begin{equation}
u_{\Sht}^{*}\Phi_{M}^{!}(\z_{1}\times\cdots\times\z_{r} \times [M])\cong \Phi_{M'}^{!}(\z'_{1}\times\cdots\times\z'_{r} \times [M])\in \Ch_{0}(\Sht'_{C}).
\end{equation}
Restricting to $\Sht'_{C}|_{s'}\cong \Sht_{C}|_s$ and taking degrees we get
\begin{equation}\label{degSS'}
\source{higher-siegel-weil-unitary-nonsingular-terms}
\deg(\Phi_{M}^{!}(\z_{1}\times\cdots\times\z_{r}  \times [M] )|_{\Sht_{C}|_s})=\deg(\Phi_{M'}^{!}(\z'_{1}\times\cdots\times\z'_{r}  \times [M])|_{\Sht'_{C}|_{s'}}).
\end{equation}
On the other hand, letting $f':M'\to S'$, by smooth base change we have
\begin{equation}\label{trSS'}
\source{higher-siegel-weil-unitary-nonsingular-terms}
\Tr(f_{!}\cla(\z_{1})\c\cdots\c f_{!}\cla(\z_{r})\c\Frob_{s}, (Rf_{!}\Qlbar)_{s}) = \Tr(f'_{!}\cla(\z'_{1})\c\cdots\c f'_{!}\cla(\z'_{r})\c\Frob_{s'}, (Rf'_{!}\Qlbar)_{s'}).
\end{equation}
Since the right sides of \eqref{degSS'} and \eqref{trSS'} are equal by the scheme case that is already proven, the left sides of \eqref{degSS'} and \eqref{trSS'} are also equal, proving the proposition in general.
\end{proof}




Recall the graded perverse sheaf $\cK^{\mrm{Int}}_{d}(T)$ on $\Herm_{2d}$ from Definition \ref{def:Kd}. 

\begin{cor}
\source{higher-siegel-weil-unitary-nonsingular-terms}
\notready\label{cor: trace of correspondence} Let $(\cQ,h_{\cQ})\in \Herm_{2d}(k)$. We have
\source{higher-siegel-weil-unitary-nonsingular-terms}
\begin{eqnarray*}
\deg[\cZ^{r}_{\cQ}]&=&\Tr( [\Hk^{1}_{\Lagr_{2d}}]^r\c\Frob, (R(\upsilon_{2d})_{*}\Qlbar)_{\cQ}) = \sum_{i=0}^d (d-2i)^r \Tr(\Frob, \Spr^{\Herm}_{2d}[\r_{i}]_{\cQ})  \\
&=& \frac{1}{(\log q)^r} \left(\frac{d}{ds}\right)^r\Big|_{s=0} \left(q^{ds}\Tr(\Frob, \cK_d^{\mrm{Int}}(q^{-2s})_{\cQ})\right).
\end{eqnarray*}
\end{cor}
%\begin{proof} For the first equality, we would like to apply Proposition \ref{p:Lef} to $S=\Herm_{2d}$, $M=\Lagr_{2d}$, $C_{i}=\Hk^{1}_{\Lagr_{2d}}$ and $\z_{i}=[\Hk^{1}_{\Lagr_{2d}}]$. The only issue is that $\Herm_{2d}$ is not a scheme. However for any stack $S$ locally of finite type over $k$ and $s\in S(k)$, there exists a scheme $S'$ of finite type over $k$ and a smooth map $u: S'\to S$ such that $u(S'(k))$ contains $s$.  Indeed, pick any smooth map $u_{1}: S_{1}\to S$ with $S_{1}$ a scheme of finite type over $k$, such that  $s$ is contained in the image of $u$. Let $s_{1}\in S_{1}(\F_{q^{m}})$ be a point that maps to $s$. Then consider the $m$-fold fibered product $(S_{1}/S)^{m}$
%equipped with the $\Frob$-descent datum given by $(x_{1},\cdots, x_{m})\mapsto (\Frob(x_{m}), \Frob(x_{1}),\cdots, \Frob(x_{m-1}))$. This gives a descent of $(S_{1}/S)^{m}$ to $S'$ over $k$ equipped with a map $u:S'\to S$ which is still smooth since $u_{1}$ is. Now $s_{1}$ gives rise to a $k$-point $s'=(s_{1},\Frob(s_{1}),\cdots, \Frob^{m-1}(s_{1}))\in S'(k)$ such that $u(s')=s$. 
%
%Apply the above remark to get a smooth map $u: S'\to S=\Herm_{2d}$ with $s'\in S(k)$ such that $u(s')=\cQ$. Apply Proposition \ref{p:base change} to the diagram \eqref{Sht Lagr} (as diagram \eqref{ShtH}) along the base change map $u$. We realize $\cZ^{r}_{\cQ}$ as the $s'$-component of $\Sht'_{H}$, in terms of the notation from \eqref{Sht'H}. Since $\Hk^{1}_{\Lagr_{2d}}$ and $\Lagr_{2d}$ are representable over $\Herm_{2d}$, all items in the diagram \eqref{Sht'H} are schemes. We then apply Proposition \ref{p:Lef} to the diagram \eqref{Sht'H}. 
\begin{proof}
The first equality is an application of Proposition \ref{p:Lef} to the case $S=\Herm_{2d}$, $M=\Lagr_{2d}$, $C_{i}=\Hk^{1}_{\Lagr_{2d}}$ and $\z_{i}=[\Hk^{1}_{\Lagr_{2d}}]$.  The second equality follows from Proposition \ref{p:action Hk1}. The third one is a direct calculation.
\end{proof}

Combining Corollary \ref{cor: trace of correspondence} with Corollary \ref{c:Zloc} we get:
\begin{cor}
\source{higher-siegel-weil-unitary-nonsingular-terms}
\notready\label{cor: KR degree as trace of Frob} Let $(\cE,a)\in \cA_{d}(k)$ with image $(\cQ,h_{\cQ})\in \Herm_{2d}(k)$. Then we have 
\source{higher-siegel-weil-unitary-nonsingular-terms}
\[
\deg [\Cal{Z}^{r}_{\Cal{E}}(a)] = \frac{1}{(\log q)^r} \left(\frac{d}{ds}\right)^r\Big|_{s=0} \left(q^{ds}\Tr(\Frob, \cK_d^{\mrm{Int}}(q^{-2s})_{\cQ})\right).
\]
\end{cor}


%
%\subsection{The action of the Hecke correspondence for $\cM_{d}$}
%Recall from Definition \ref{defn: hitchin hecke} that $\cH_{d}:=\Hk^{1}_{\cM_{d}}$ classifies diagrams of vector bundles 
%\[
%\begin{tikzcd} & \Cal{F} \ar[dd, dashed, "+x'-\s(x')"] \\ \Cal{E} \ar[ur] \ar[dr] \\ & \Cal{F}' \end{tikzcd} 
%\]
%where dotted arrow denotes an upper modification of the Hermitian bundle $(\cF,h)$ at $x'$ and lower modification at $\s(x')$ to get $(\cF',h')$. Here $(\cE\to\cF, h)\in \cM_{d}$ and $(\cE\to \cF',h')\in \cM_{d}$. We view $\cH_{d}$ as a self-correspondence $\cM_{d}$ over $\cA_{d}$:
%\begin{equation}\label{corr M}
%\begin{tikzcd}
%& \cH_{d} \ar[dl] \ar[dr] \\
%\Cal{M}  \ar[dr, "f"] & & \Cal{M} \ar[dl, "f"] \\
%& \Cal{A}
%\end{tikzcd}
%\end{equation}
%The goal of this subsection is to calculate the action of $\cH_{d}$ on the direct image complex $Rf_{d*}\Qlbar$. 
%
%We have a similar self-correspondence $C_{d}$ of $\Lagr_{2d}$ over $\Herm_{2d}$: it classifies $(\cQ,h,\cL,\cL')$ where $\cL\subset \cQ$ and $\cL'\subset \cQ'$ are Lagrangians such that (over geometric points of any test scheme) $\cL/(\cL\cap \cL')$ has length one. 
%\begin{equation}\label{corr Herm}
%\xymatrix{& C_{d} \ar[dl]\ar[dr] \\
%\Lagr_{2d}\ar[dr]_{\ph} && \Lagr_{2d}\ar[dl]^{\ph}\\
%& \Herm_{2d}}
%\end{equation}
%We have a canonical map $g_{\cH}: \cH_{d}\to C_{d}$ sending $(\cE\to\cF, h, \cE\to cF', h')$ to $(\s^{*}\cE^{\vee}/\cE, \ov h, \cF/\cE, \cF'/\cE)$.   The maps $g_{\cH}, g_{\cM}$ and $g$ exhibit the diagram \eqref{corr M} as the pullback of the diagram \eqref{corr Herm} via the base change $g: \cA_{d}\to\Herm_{2d}$. 
% 
%
%%where \mnote{this notation should be explained somewhere earlier}
%%\[
%%C = \left\{ \begin{tikzcd} & \Cal{F} \ar[dd, dashed, "+x'-x''"] \\ \Cal{E} \ar[ur] \ar[dr] \\ & \Cal{F}' \end{tikzcd} \right\}
%%\]
%%and calculate the action of $[C]$ on $Rf_* \ol{\Q}_{\ell}$, because by \S \ref{ssec: geometric reform} we have 
%%\[
%%\Tr(\Frob \circ [C]^r, (Rf_! \ol{\Q}_{\ell})_a) = \deg \Sht_{\Cal{M}}^r(a).
%%\]
%
%By Corollary \ref{c:Lag decomp} and $Rf_{*}\Qlbar\cong g^{*}R\ph_{*}\Qlbar$, we get a decomposition
%\begin{equation}\label{KAi}
%Rf_*\Qlbar\cong \bigoplus_{i=0}^d K_{\cA, i}
%\end{equation}
%where $K_{\cA,i}=g^{*}\Spr^{\Herm}_{2d}[\r_{i}]$ is shifted perverse since $g$ is smooth.
%
%\begin{prop} The action of the cohomological correspondence $[\cH_{d}]$ on $Rf_{*}\Qlbar$ preserves the decomposition \eqref{KAi}, and it acts on $K_{\cA,i}$ by multiplication by $(d-2i)$.
%\end{prop}
%\begin{proof}
%By Proposition \ref{prop:Ad Herm smooth },  $g$ is smooth. Hence the cohomological correspondence $[\cH_{d}]$ on the constant sheaf $\Qlbar$ on $\cM_{d}$ is obtained from the cohomological correspondence $[C_{d}]$ on the constant sheaf $\Qlbar$ on $\Lagr_{2d}$ by a smooth pullback. The statement about the $[\cH_{d}]$ action on $Rf_{*}\Qlbar$ then follows from the similar statement for the $[C_{d}]$ action on $R\ph_{*}\Qlbar=\op  K_{i}$ where $K_{i}=\Spr^{\Herm}_{2d}[\r_{i}]$. Since $K_{i}$ are full support perverse sheaves, it suffices to prove the same statement on the open dense substack $\Herm_{2d}^{\c}$, the preimage of $X^{\c}_{d}$. Now over $\Herm_{2d}^{\c}$, the map $\Herm_{2d}^{\c}\to X^{\c}_{d}$ is smooth, and $C_{d}$ is the pullback of the incidence correspondence $I'_{d}$ over $X'_{d}$ (see \cite[proof of Proposition 8.3]{YZ}). This in turn reduces to checking the statement for the action of $[I'_{d}]$ on the direct image of $\nu_{d}: X'_{d}\to X_{d}$, which is done in the proof of \cite[Proposition 8.3]{YZ}.
%\end{proof}

%The smallness of $f$ implies that the $K_i'$ are perverse. The correspondence $[C]$ acts on $K_i$ by multiplication by $(2d-i)$. \mnote{is this right?} 


%The second equality above uses that $r$ is even.


%\begin{proof}
%The first equality follows from Theorem \ref{}
%Since $K_{\cA,i}$ is the pullback of $\Spr^{\Herm}_{2d}[\r_{i}]$ via $g: \cA_{d}\to \Herm_{2d}$, we may restate Corollary \ref{cor: trace of correspondence} as
%\begin{equation}
%\Tr(\Frob \circ [\cH_{d}]^r, (Rf_*\ol{\Q}_{\ell})_{\ov a})  =\sum_{i=0}^d (d-2i)^r \Tr(\Frob, \Spr^{\Herm}_{2d}[\r_{i}]_{\cQ}). 
%\end{equation}
%The right side above is easily seen to be equal to
%\begin{equation}
%\frac{1}{(\log q)^r} \left(\frac{d}{ds}\right)^r\mid_{s=0} \left(q^{ds}\Tr(\Frob_q, \cK_{d}(q^{-2s})_{\cQ})\right).
%\end{equation}
%On the other hand,  the left side is equal to the degree of $\cZ^{\mu}_{\cE}(a)$ by \footnote{refer to previous section...}
%\end{proof}
%%% added 12.30 %%%%
\subsection{Symmetry} This subsection is not used in the proof of the main theorem. The graded perverse sheaf $\cK^{\Int}_{d}(T)$ has a palindromic symmetry that we spell out. First, the \'etale double covering $\nu:X'\to X$ gives a local system $\y_{X'/X}$ on $X$ with monodromy in $\pm1$.  It induces a local system $\y_{d}$ on $X_{d}$ with monodromy in $\pm1$: its stalk at a divisor $x_{1}+\cdots+x_{d}\in X_{d}(\ov{k})$ is $\ot_{i=1}^{d}(\y_{X'/X})_{x_{i}}$. Let 
\begin{equation}
\y^{\Herm}_{2d}:=s^{\Herm*}_{2d}\y_{d},
\end{equation}
where $s^{\Herm}_{2d}: \Herm_{2d}\to X_{d}$ is the support map. This is a rank-one local system on $\Herm_{2d}$ with monodromy in $\pm1$.

\begin{lemma}
\source{higher-siegel-weil-unitary-nonsingular-terms}
\notready We have a canonical isomorphism of perverse sheaves on $\Herm_{2d}$:
\begin{equation}
\Spr^{\Herm}_{2d}[\r_{d}]\cong\y^{\Herm}_{2d}.
\end{equation}
\end{lemma}
\begin{proof}
By Proposition \ref{p:SH Spr}(2), $\Spr^{\Herm}_{2d}[\r_{d}]$ is the middle extension of its restriction to the open dense substack $\Herm_{2d}^{\c}$ (preimage of $X^{\c}_{d}$). The same is true for $\y_{d}$  because it is a local system and $\Herm_{2d}$ is smooth. Therefore it suffices to check their equality over $\Herm_{2d}^{\c}$, over which both are obtained by pushing out the $W_{d}$-torsor $(X'^{d})^{\c}\to X_{d}^{\c}$ along the character $\chi_{d}: W_{d}\to \{\pm1\}$.
\end{proof}

\begin{lemma}
\source{higher-siegel-weil-unitary-nonsingular-terms}
\notready There is an isomorphism of graded perverse sheaves on $\Herm_{2d}$
\begin{equation}
T^{d}\cK_d^{\Int}(T^{-1})\cong \y^{\Herm}_{2d}\ot \cK^{\Int}_{d}(T).
\end{equation}
\end{lemma}
\begin{proof}
The equality amounts to 
\begin{equation}
\Spr^{\Herm}_{2d}[\r_{d-i}]\cong \y^{\Herm}_{2d}\ot\Spr^{\Herm}_{2d}[\r_{i}].
\end{equation}
Both sides are middle extensions from $\Herm^{\c}_{2d}$ by Proposition \ref{p:SH Spr}(2), over which they correspond to representations $\r_{d-i}$ and $\chi_{d}\ot \r_{i}$ of $W_{d}$. By definition, 
\begin{equation}
\chi_{d}\ot\r_{i}\cong \chi_{d}\ot\Ind^{W_{d}}_{W_{i}\times W_{d-i}}(\chi_{i}\bt\one).
\end{equation}
Inserting $\chi_{d}|_{W_{i}\times W_{d-i}}\cong \chi_{i}\bt \chi_{d-i}$ to the right side above gives 
\begin{equation}
\Ind^{W_{d}}_{W_{i}\times W_{d-i}}(\chi_{d}|_{W_{i}\times W_{d-i}}\ot(\chi_{i}\bt\one))\cong \Ind^{W_{d}}_{W_{i}\times W_{d-i}}(\one\bt\chi_{d-i})\cong \r_{d-i}.
\end{equation}
\end{proof}

\begin{lemma}
\source{higher-siegel-weil-unitary-nonsingular-terms}
\notready
If $(\cQ,h_{\cQ})\in \Herm_{2d}(k)$ is the image of some $(\cE,a)\in \cA_{d}(k)$, then $\Tr(\Frob, \y^{\Herm}_{2d}|_{\cQ})=1$.
\end{lemma}
\begin{proof} If $X'/X$ is split, then $\y^{\Herm}_{2d}$ is trivial, and there is nothing to prove. Below we assume $X'/X$ is nonsplit. The local system $\y_{d}$ on $X_{d}$ is pulled back from a local system $\y_{\Pic}$ on $\Pic_{X}$ via the Abel-Jacobi map $\mrm{AJ}_{d}:X_{d}\to \Pic^{d}_{X}\subset \Pic_{X}$. The Frobenius trace function of $\y_{\Pic}$ is the id\`ele class character 
\[
\y_{F'/F}: F^{\times}\bs\BA^{\times}_{F}/\wh\cO^{\times}=\Pic_{X}(k)\to\{\pm1\}
\]
trivial on the image of $\Nm_{X'/X}: \Pic_{X'}(k)\to \Pic_{X}(k)$. Denote by $\det_{X}(\cQ)$ the image of $\cQ$ under $\Herm_{2d}\to X_{d}\xr{\mrm{AJ}_{d}} \Pic^{d}_{X}$. We have $\y^{\Herm}_{2d}|_{\cQ}\cong \y_{\Pic}|_{\det_{X}(\cQ)}$ as $\Frob$-modules. Now $(\cQ,h_{\cQ})$ comes from $(\cE,a)$, which implies
\begin{equation*}
{\det}_{X}(\cQ)\cong \Nm_{X'/X}(\det\cE)^{-1}\ot \om_{X}^{\ot n}.
\end{equation*}
By \cite[p.291, Theorem 13]{Weil}, $\om_{X}$ is a square in $\Pic_{X}(k)$, hence $\y_{F'/F}(\om_{X}^{\ot n})=1$. Since $\y_{F'/F}(\Nm_{X'/X}(\det\cE))=1$, we see that $\y_{F'/F}(\det_{X}(\cQ))=1$, hence $\Tr(\Frob, \y^{\Herm}_{2d}|_{\cQ})=1$.
\end{proof}


\begin{cor}
\source{higher-siegel-weil-unitary-nonsingular-terms}
\notready\label{cor:FE Int}
\source{higher-siegel-weil-unitary-nonsingular-terms}
Let $(\cE,a)\in \cA_{d}(k)$ with image  $(\cQ,h_{\cQ})\in \Herm_{2d}(k)$. Then $s\mapsto q^{ds}\Tr(\Frob, \cK_d^{\mrm{Int}}(q^{-2s})_{\cQ})$ is an even function in $s$. In particular, its odd order derivatives at $s=0$ vanish.
\end{cor}
By Corollary \ref{cor: KR degree as trace of Frob}, this implies $\deg[\cZ^{r}_{\cE}(a)]=0$ for $r$ odd. However, we know from Lemma \ref{lem: shtuka empty parity} that $\Sht^{r}_{U(n)}=\vn$ when $r$ is odd, which implies $\cZ^{r}_{\cE}(a)=\vn$. 





\part{The comparison}
